% ---------------------------------------------------------------
% Baybayani Technical Evaluation Instrument
% Product quality, based on ISO/IEC 25010:2023
% Respondents: technical stakeholders (developers, testers and
% system architects).
% Build: pdflatex technical-instrument-iso25010.tex (run twice)
% ---------------------------------------------------------------
\documentclass[10pt]{article}
% ---------------------------------------------------------------
% Study details shared by both evaluation instruments.
% Edit here once; both instruments pick the changes up on rebuild.
% Source: VSU appointment letters for Project DCST.0824-0825
% (September 10, 2024).
% ---------------------------------------------------------------

\newcommand{\SystemName}{Baybayani}
\newcommand{\ProjectTitle}{Farm to Table: A Management Information System for
  Farmers and Food Business Establishments in Baybay City, Leyte}
\newcommand{\ProjectCode}{DCST.0824-0825}

\newcommand{\ProjectLeader}{Mr.~Nelson H.~Tejara}
% Study/Component Leaders, alphabetical by surname
\newcommand{\ComponentLeaders}{%
  Mr.~Jomari Joseph A.~Barrera,
  Dr.~Jonah Flor O.~Maaghop,
  Mr.~Rodney M.~Maniego Jr.,
  Dr.~Jude B.~Rola, and
  Dr.~Magdalene C.~Unajan}

\newcommand{\University}{Visayas State University}
\newcommand{\Faculty}{Faculty of Computing}
\newcommand{\Department}{Department of Computer Science and Technology}
\newcommand{\Address}{Visca, Baybay City, Leyte 6521}

\usepackage{instrument}

\renewcommand{\InstrumentShortName}{Technical Evaluation (ISO/IEC 25010:2023)}

\begin{document}

\InstrumentTitle{Technical Evaluation Instrument}%
  {Product Quality of the \SystemName{} System, based on ISO/IEC 25010:2023\\
   For developers, testers and system architects}

% ===============================================================
\PartHeading{Part I. Informed Consent and Data Privacy Notice}

Faculty researchers of the \Department{} of \University{} are evaluating the
product quality of \SystemName, the farm-to-table ordering system of the research
project named above. We are asking you, as a technical stakeholder, to rate the
system against the nine product quality characteristics of ISO/IEC 25010:2023,
based on your own review, testing or analysis.

\begin{itemize}
  \item \textbf{Voluntary.} Taking part is voluntary. You may skip any item or stop at
    any time.
  \item \textbf{Time.} Rating takes about 30 to 45 minutes after you have reviewed or
    tested the system.
  \item \textbf{Confidentiality.} Your name appears only on this page, which is filed
    separately from your ratings. Ratings are identified only by a respondent number
    and are reported as group summaries.
  \item \textbf{Data privacy.} Your information is collected, stored and used only for
    this evaluation, in line with the Data Privacy Act of 2012 (Republic Act
    No.~10173). Only the researchers can access it, and it is disposed of securely
    once the project is completed.
  \item \textbf{Confidential material.} Access to the source code, builds and test
    accounts is given only for this evaluation. Please do not copy, share or keep
    them after the evaluation.
\end{itemize}

For questions about this evaluation, approach any of the researchers at the
\Department, \University.

\vspace{4pt}
\opt{} I have read and understood this notice, and I agree to take part.

\vspace{0.9cm}
\noindent
\begin{tabular}{@{}p{8.6cm}@{\hspace{1cm}}p{5cm}@{}}
  \rule{\linewidth}{0.4pt} & \rule{\linewidth}{0.4pt}\\
  \small Signature over printed name & \small Date
\end{tabular}

\clearpage
% ===============================================================
\PartHeading{Part II. Evaluator Profile}

Tick (\checkmark) the box that applies.

{\small
\renewcommand{\arraystretch}{1.55}
\begin{tabularx}{\linewidth}{|>{\raggedright\arraybackslash}p{5.1cm}|X|}
\hline
\textbf{1. Primary role} &
  \opt{Software developer} \opt{Software tester / QA} \opt{System architect}
  \opt{Other: \blank[3cm]}\\ \hline
\textbf{2. Affiliation} &
  \opt{Academe (faculty)} \opt{Industry} \opt{Government}
  \opt{Other: \blank[3cm]}\\ \hline
\textbf{3. Years of professional experience in the role} &
  \opt{Less than 1} \opt{1--3} \opt{4--6} \opt{7--10} \opt{More than 10}\\ \hline
\textbf{4. Involvement in developing \SystemName} &
  \opt{None} \opt{Consulted or reviewed during development}
  \opt{Member of the development team}\\ \hline
\textbf{5. Familiarity with the technologies used}\newline
  {\footnotesize (N = none, B = basic, P = proficient)} &
  \renewcommand{\arraystretch}{1.2}%
  \begin{tabular}[t]{@{}l@{\hspace{1.2em}}l@{}}
  React and TypeScript                  & \opt{N} \opt{B} \opt{P}\\
  Supabase / PostgreSQL (RLS, Realtime) & \opt{N} \opt{B} \opt{P}\\
  Capacitor / Android                   & \opt{N} \opt{B} \opt{P}\\
  Web security and testing practice     & \opt{N} \opt{B} \opt{P}\\
  \end{tabular}\\ \hline
\textbf{6. How you evaluated the system}\newline
  {\footnotesize (tick all that apply)} &
  \opt{Hands-on use of the web app} \opt{Hands-on use of the Android app}
  \opt{Source code review} \opt{Database schema and access-policy review}
  \opt{Functional or exploratory testing} \opt{Performance or load testing}
  \opt{Security testing} \opt{Architecture or documentation review}\\ \hline
\end{tabularx}}

\vspace{6pt}
\textbf{Evaluation environment} (filled in by the evaluator or the researchers)

{\small
\renewcommand{\arraystretch}{1.55}
\begin{tabularx}{\linewidth}{|>{\raggedright\arraybackslash}p{5.1cm}|X|}
\hline
Version evaluated (Git commit or release) & \\ \hline
Web URL tested & \\ \hline
Android device and version tested & \\ \hline
Browsers and devices tested & \\ \hline
Test accounts used & \opt{Customer} \opt{Administrator}\\ \hline
Evaluation date(s) & \\ \hline
\end{tabularx}}

% ===============================================================
\PartHeading{Part III. Product Quality Evaluation}

\textbf{Directions.} Rate each statement according to the evidence you gathered
from your review or testing. Shade or tick (\checkmark) one circle per item. Where a
rating is based on a specific observation (a test case, a file, a screen, a log),
note it in the observation box below the characteristic, especially for ratings of
2 or lower.

\begin{center}
\begin{minipage}{\linewidth}
\centering
\small
\renewcommand{\arraystretch}{1.3}
\begin{tabular}{|c|c|l|c|}
\hline
\rowcolor{headshade}
\textbf{Rating} & \textbf{Code} & \textbf{Description} & \textbf{Mean range}\\ \hline
5 & SA & Strongly Agree    & 4.21--5.00\\ \hline
4 & A  & Agree             & 3.41--4.20\\ \hline
3 & N  & Neutral           & 2.61--3.40\\ \hline
2 & D  & Disagree          & 1.81--2.60\\ \hline
1 & SD & Strongly Disagree & 1.00--1.80\\ \hline
\end{tabular}\\[2pt]
{\footnotesize The mean range is used by the researchers to interpret the results.}
\end{minipage}
\end{center}

\StartNumbering{}

% ---------------------------------------------------------------
\Characteristic{1. Functional Suitability}{The degree to which the system provides
  functions that meet the stated and implied needs of its users.}
\begin{ratingtable}
\SubChar{1.1 Functional completeness}
\EnItem{The customer functions cover the intended tasks: sign-up with valid-ID upload,
  browsing and search, cart, checkout, order tracking and cancellation, chat,
  announcements and profile settings.}
\EnItem{The administrator functions cover the intended tasks: products and variants,
  stock movements, order processing (including walk-in orders), user approval,
  announcements, store hours, dashboard and report export.}
\SubChar{1.2 Functional correctness}
\EnItem{Computed values are correct: cart and order totals, retail and wholesale
  pricing, and dashboard revenue and order figures.}
\EnItem{Stock levels stay consistent with recorded movements after sales, losses,
  acquisitions, adjustments and cancellation restocks.}
\EnItem{Orders keep the variant and price at the time of purchase even when the product
  is later edited or its price changes.}
\SubChar{1.3 Functional appropriateness}
\EnItem{Each function achieves its task without unnecessary steps or screens.}
\end{ratingtable}
\CommentBox[2]{Observations and evidence:}

% ---------------------------------------------------------------
\Characteristic{2. Performance Efficiency}{The degree to which the system performs
  its functions within specified time and throughput, and uses resources
  efficiently, under the stated conditions.}
\begin{ratingtable}
\SubChar{2.1 Time behaviour}
\EnItem{Pages, product searches and order actions respond within acceptable time on
  typical mobile and Wi-Fi connections.}
\EnItem{Real-time updates (chat messages, order status, notifications) appear without
  noticeable delay.}
\SubChar{2.2 Resource utilization}
\EnItem{The apps use device memory, battery, storage and mobile data reasonably (e.g.,
  compressed images, lazy-loaded pages, paginated lists).}
\EnItem{Database queries and realtime subscriptions fetch only the data a screen
  needs.}
\SubChar{2.3 Capacity}
\EnItem{The system handles the expected number of products, users and concurrent
  orders (e.g., peak ordering before closing time) without degradation.}
\end{ratingtable}
\CommentBox[2]{Observations and evidence:}

% ---------------------------------------------------------------
\Characteristic{3. Compatibility}{The degree to which the system can exchange
  information with other systems and perform its functions while sharing the same
  environment and resources.}
\begin{ratingtable}
\SubChar{3.1 Co-existence}
\EnItem{The Android app and the web app run alongside other applications on the same
  device without conflicts or adverse effects.}
\SubChar{3.2 Interoperability}
\EnItem{The system exchanges data correctly with its external services: Supabase
  Auth, Storage and Realtime, and Firebase Cloud Messaging.}
\EnItem{Reports exported to Excel open correctly and contain the expected data.}
\EnItem{The web and Android clients show the same data for the same account at the
  same time.}
\end{ratingtable}
\CommentBox[2]{Observations and evidence:}

% ---------------------------------------------------------------
\Characteristic{4. Interaction Capability}{The degree to which the intended users
  can interact with the system to complete their tasks, through its user interface.}
\begin{ratingtable}
\SubChar{4.1 Appropriateness recognizability}
\EnItem{Users can tell from the interface whether the system fits their needs (e.g.,
  the shop, cart and admin sections are clear from first use).}
\SubChar{4.2 Learnability}
\EnItem{New customers and administrators can learn the main tasks in a short time
  without formal training.}
\SubChar{4.3 Operability}
\EnItem{Controls are easy to operate on both the mobile and desktop layouts.}
\SubChar{4.4 User error protection}
\EnItem{Input validation, confirmation dialogs, unsaved-changes warnings and
  insufficient-stock checks protect users from making errors.}
\SubChar{4.5 User engagement}
\EnItem{The interface is attractive and encourages continued use.}
\SubChar{4.6 Inclusivity}
\EnItem{The system can be used by people of different ages, backgrounds and levels of
  digital experience, including farmers and small food businesses.}
\SubChar{4.7 User assistance}
\EnItem{The system supports users with different abilities (e.g., readable text
  sizes, colour contrast, light and dark themes, labels for screen readers, adequate
  touch targets).}
\SubChar{4.8 Self-descriptiveness}
\EnItem{Status labels, empty states, error messages and notifications make the
  system's state and next steps clear without a manual.}
\end{ratingtable}
\CommentBox[2]{Observations and evidence:}

% ---------------------------------------------------------------
\Characteristic{5. Reliability}{The degree to which the system performs its
  functions without failure under specified conditions for a specified period.}
\begin{ratingtable}
\SubChar{5.1 Faultlessness}
\EnItem{The system performs its functions without errors or crashes under normal
  operation.}
\SubChar{5.2 Availability}
\EnItem{The system is up and reachable whenever customers and administrators need it.}
\SubChar{5.3 Fault tolerance}
\EnItem{The system keeps working, with clear messages, when the network drops or a
  request to the backend fails; no partial orders are created.}
\SubChar{5.4 Recoverability}
\EnItem{After an interruption, the system restores the user's session, cart and data
  (e.g., the server-side cart, re-login, pull-to-refresh).}
\EnItem{Data can be restored from backups if the database is lost or corrupted.}
\end{ratingtable}
\CommentBox[2]{Observations and evidence:}

% ---------------------------------------------------------------
\Characteristic{6. Security}{The degree to which the system protects information and
  data so that people and other systems have access appropriate to their
  authorization, and defends against attack.}
\begin{ratingtable}
\SubChar{6.1 Confidentiality}
\EnItem{Users can access only data they are authorized to see (e.g., customers see
  only their own orders, cart and chats; valid-ID images are not publicly
  accessible).}
\EnItem{Row-level security policies are enabled on every table that holds user or
  business data.}
\SubChar{6.2 Integrity}
\EnItem{Prices, stock, order status and user roles cannot be changed by unauthorized
  users, including through direct API calls with the public key.}
\SubChar{6.3 Non-repudiation}
\EnItem{Orders, stock movements and status changes are recorded with timestamps so that
  they can be proven to have taken place.}
\SubChar{6.4 Accountability}
\EnItem{Administrative actions (stock adjustments, order status changes, user
  approvals) can be traced to the account that performed them.}
\SubChar{6.5 Authenticity}
\EnItem{User identity is verified through authentication, email verification, password
  reset and administrator approval of new accounts.}
\SubChar{6.6 Resistance}
\EnItem{The system resists common attacks (e.g., injection, cross-site scripting,
  brute-force login, exposure of secret keys in the client or repository).}
\end{ratingtable}
\CommentBox[2]{Observations and evidence:}

% ---------------------------------------------------------------
\Characteristic{7. Maintainability}{The degree to which the system can be modified
  effectively and efficiently by its intended maintainers.}
\begin{ratingtable}
\SubChar{7.1 Modularity}
\EnItem{The code is organized into separate parts (pages, components, data hooks,
  models, route guards) so that a change to one has little effect on others.}
\SubChar{7.2 Reusability}
\EnItem{Components, hooks and utilities are reused rather than duplicated.}
\SubChar{7.3 Analysability}
\EnItem{The causes of defects and the effects of a planned change can be found easily
  (readable code, TypeScript types, linting, README documentation).}
\SubChar{7.4 Modifiability}
\EnItem{Features can be changed or added without introducing defects or lowering
  quality.}
\SubChar{7.5 Testability}
\EnItem{Test criteria can be set and tests run effectively (e.g., automated tests,
  test data, separate test environment).}
\end{ratingtable}
\CommentBox[2]{Observations and evidence:}

% ---------------------------------------------------------------
\Characteristic{8. Flexibility}{The degree to which the system can be adapted to
  changes in its requirements, contexts of use or system environment.}
\begin{ratingtable}
\SubChar{8.1 Adaptability}
\EnItem{The system works across supported browsers, Android versions and screen
  sizes.}
\SubChar{8.2 Scalability}
\EnItem{The system can handle growth in products, users, orders or additional sellers
  by scaling its resources.}
\SubChar{8.3 Installability}
\EnItem{The Android APK installs, updates and uninstalls cleanly, and the web app can
  be built and deployed by following the README.}
\SubChar{8.4 Replaceability}
\EnItem{The system can replace the current manual, paper or social-media ordering
  process for the same purpose.}
\end{ratingtable}
\CommentBox[2]{Observations and evidence:}

% ---------------------------------------------------------------
\Characteristic{9. Safety}{The degree to which the system, under defined conditions,
  avoids states that put human life, health, property or the environment at risk.
  For this system, the items address harm to the farmer group's stock, orders,
  money and records.}
\begin{ratingtable}
\SubChar{9.1 Operational constraint}
\EnItem{The system keeps operations within safe limits (e.g., no ordering beyond
  available stock, no negative stock, no ordering outside store hours).}
\SubChar{9.2 Risk identification}
\EnItem{The system identifies conditions that could cause loss (e.g., low-stock
  alerts, insufficient-stock detection, orders affected by closing time).}
\SubChar{9.3 Fail safe}
\EnItem{When an operation fails, the system leaves data in a safe state (e.g., a failed
  checkout neither creates an order nor deducts stock).}
\SubChar{9.4 Hazard warning}
\EnItem{The system warns users before actions that are hard to undo (e.g., deleting a
  product or variant, cancelling an order, adjusting stock, closing with unsaved
  changes).}
\SubChar{9.5 Safe integration}
\EnItem{Integrated services (push notifications, edge functions, Excel export) do not
  expose data or leave the system in an inconsistent state.}
\end{ratingtable}
\CommentBox[2]{Observations and evidence:}

% ===============================================================
\par\noindent\begin{minipage}{\linewidth}
\PartHeading{Part IV. Overall Assessment}

\CommentBox[3]{1. What are the strengths of \SystemName{} as a product?}
\vspace{4pt}
\CommentBox[3]{2. What defects or risks did you find? Give the steps to reproduce,
  the file or the screen where possible.}
\vspace{4pt}
\CommentBox[3]{3. What do you recommend before wider deployment?}

\vspace{10pt}
\begin{center}
\textbf{Thank you for your time and expert assessment.}
\end{center}
\end{minipage}

\end{document}
